
$X=\Spec(B)$ ; on a le diagramme commutatif :
\[
\begin{tikzcd}[column sep=large,row sep=large]
& A[T_1,\ldots,T_n] \arrow[dl,"v"']\\
B & A \arrow[l,"u"'] \arrow[u]
\end{tikzcd}
\]
où les morphismes $u$ et $v$ sont surjectifs. Soient $t_i$, pour $i=1,\ldots,n$, des éléments de $A$ tels que pour tout $i$, on ait
\[
u(t_i)=v(T_i).
\]
On définit un homomorphisme $w:A[T_1,\ldots,T_n]\longrightarrow A$ en posant $w(T_i)=t_i$. On obtient ainsi une immersion fermée $s:V\to V'$ qui est une section de $p$, et qui est telle que $j=s\circ i$. De plus, étant une section d'un morphisme lisse, c'est une immersion régulière (VII 1.10). Sachant que $s$ et $s\circ i$ sont régulières, on en déduira que $i$ est régulière si on montre que la suite :
\[
0\longrightarrow i^*(\underline{K}/\underline{K}^2)\longrightarrow \underline{J}/\underline{J}^2\longrightarrow \underline{I}/\underline{I}^2\longrightarrow0
\]
est exacte et localement scindée, $\underline{I}$, $\underline{J}$ et $\underline{K}$ désignant respectivement les idéaux définissant $i$, $j$, et $s$ (VII 1.7). Il suffit pour cela de montrer que $i^*(\underline{K}/\underline{K}^2)$ est facteur direct de $\underline{J}/\underline{J}^2$. Or on voit tout de suite que $j'$ se déduit de $s$ par le changement de base $i':X'\to V'$. En outre, les immersions $j'$ et $s$ sont régulières (VII 1.10). Par conséquent (VII 1.2), l'idéal qui définit $j'$ est $i'^*(\underline{K})$, et le faisceau conormal de $j'$ est $i^*(\underline{K}/\underline{K}^2)$. La suite exacte des faisceaux conormaux relative à la décomposition $j=i'\circ j'$ donne donc en particulier un homomorphisme : $\underline{J}/\underline{J}^2\longrightarrow i^*(\underline{K}/\underline{K}^2)$. Il est immédiat de vérifier que le composé : $i^*(\underline{K}/\underline{K}^2)\longrightarrow \underline{J}/\underline{J}^2\longrightarrow i^*(\underline{K}/\underline{K}^2)$ est l'identité (du point de vue affine, cela correspond à la décomposition de $J$ en somme directe d'un idéal de $A$, et de l'idéal engendré par les $T_i-t_i$). On a donc bien ainsi le scindage cherché.

